\chapter{Platforms and Innovation}\label{ch:b3:platforms-and-innovation}

Two features of modern markets strain the models of the earlier chapters. Some
goods are worth more the more people use them: a telephone network, a payment card,
a ride platform. And much of what firms sell is the result of research, whose
reward depends on how long others can be kept from copying it. This chapter builds
the models of \omterm{def:b3:platforms-and-innovation:network}{network effects}, of platforms that serve two groups at once, of
\omterm{def:b3:platforms-and-innovation:switching}{switching costs} and of \omterm{def:b3:platforms-and-innovation:patent}{patents}, and ends with the evidence that much of an economy's
productivity growth comes from firms entering and leaving, not from firms improving.

\begin{openingchart}{Research and development spending as a share of GDP,
1996--2023, six economies (India to 2020).}{World Bank, World Development Indicators,
GB.XPD.RSDV.GD.ZS}
\begin{tikzpicture}
\begin{axis}[omtime, width=0.7\linewidth, height=0.42\linewidth, xmin=1996, xmax=2023,
  ymin=0, ymax=7, xtick distance=5, ylabel={Percent of GDP}, legend pos=north west,
  legend style={font=\footnotesize}, legend columns=2]
  \addplot[omDef, thick] table[x=year, y=ISR] {figdata/bachelor-3/29-platforms-and-innovation/out/rnd.csv};
  \addplot[omThm, thick] table[x=year, y=KOR] {figdata/bachelor-3/29-platforms-and-innovation/out/rnd.csv};
  \addplot[black, thick] table[x=year, y=USA] {figdata/bachelor-3/29-platforms-and-innovation/out/rnd.csv};
  \addplot[omDef, dashed, thick] table[x=year, y=DEU] {figdata/bachelor-3/29-platforms-and-innovation/out/rnd.csv};
  \addplot[omThm, dashed, thick] table[x=year, y=CHN] {figdata/bachelor-3/29-platforms-and-innovation/out/rnd.csv};
  \addplot[black, dotted, thick] table[x=year, y=IND] {figdata/bachelor-3/29-platforms-and-innovation/out/rnd.csv};
  \legend{Israel, Korea, United States, Germany, China, India}
\end{axis}
\end{tikzpicture}
\end{openingchart}

% figdata: rnd_points
Israel spent 2.5~\% of GDP on research in 1996 and 6.3~\% in 2023; Korea went from 2.1 to 4.9~\%, China
from 0.6 to 2.6~\%. The United States rose from 2.5 to 3.4~\% and Germany from 2.1 to 3.2~\%, while India stayed near
0.6~\%. (The United Kingdom is left out: its series jumps by more than half a point
between 2013 and 2014, a change of measurement.) What economies buy with this
spending, and how its rewards are protected, are the subject of half this chapter.

\begin{recall}
Growth through new ideas and the replacement of incumbents are the Year 2 volume's;
the deadweight loss of monopoly is the Year 1 volume's, and dynamic games the Year 2
volume's. Panel regressions are \cref{ch:b3:panel-data}'s.
\end{recall}

\begin{notation}[Symbols of this chapter]\label{not:b3:platforms-and-innovation:symbols}
$n$ is a share or number of participants; $\alpha_i$ the value a member of side $i$ of a platform
places on each member of the other side; $f_i$ the platform's cost per member of side $i$;
$s$ a \omterm{def:b3:platforms-and-innovation:switching}{switching cost} and $\beta$ a discount factor; $s_{it}$ firm $i$'s share of employment in year $t$, used in
the decomposition only.
\end{notation}

\section{Network effects and tipping}\label{sec:b3:platforms-and-innovation:network}

\begin{definition}[Network effect; critical mass; tipping]\label{def:b3:platforms-and-innovation:network}
\index{network effect}\index{critical mass}\index{tipping}
A good has a \emph{network effect} when its value to a user rises with the number of
users --- directly, as with a telephone, or indirectly through complementary
goods, as with the software written for a computer. The \emph{critical mass} is the
smallest number of users beyond which use grows on its own. A market \emph{tips} when
small early advantages make one network take most of it.
\end{definition}

\begin{proposition}[Network demand has several equilibria]\label{prop:b3:platforms-and-innovation:rohlfs}
Let consumers' tastes $\theta$ be uniform on $[0, 1]$, and let a consumer value the service at $V\theta n$ when a share $n$ of
consumers uses it. At price $p$, a share $n > 0$ is an equilibrium if and only if $Vn(1 - n) = p$. For $0 < p < V/4$ there are two
such shares and the share 0 is also an equilibrium; the lower one is the \omterm{def:b3:platforms-and-innovation:network}{critical mass}.
\end{proposition}

\begin{proof}
If a share $n$ uses the service, a consumer joins iff $V\theta n \ge p$; the joiners are those with $\theta \ge p/(Vn)$, a
share $1 - p/(Vn)$. Consistency requires $n = 1 - p/(Vn)$, that is $Vn(1 - n) = p$, whose left side peaks at $V/4$ when $n = \tfrac12$; below
that peak the equation has two roots. With no users, no one gains by joining, so 0 is
an equilibrium. Between the roots $Vn(1 - n) > p$: more consumers want to join than use it, and
use grows towards the upper root; below the lower root use shrinks to zero.
\end{proof}

\begin{omfigure}
\begin{tikzpicture}[econ, x=7cm, y=1.2cm]
  \draw[econ axes] (0,0) -- (1.08,0) node[right, font=\scriptsize] {$n$};
  \draw[econ axes] (0,0) -- (0,2.9) node[above, font=\scriptsize] {willingness to pay};
  \draw[curve, omDef, domain=0:1, samples=60] plot (\x, {10*\x*(1-\x)});
  \draw[curve, omThm] (0,1.6) -- (1,1.6) node[right, font=\scriptsize] {price};
  \draw[guide] (0.2,0) -- (0.2,1.6);
  \draw[guide] (0.8,0) -- (0.8,1.6);
  \node[below, font=\scriptsize] at (0.2,0) {critical mass};
  \node[below, font=\scriptsize] at (0.8,0) {stable};
  \draw[->, thick] (0.35,2.05) -- (0.6,2.05);
  \draw[<-, thick] (0.03,0.7) -- (0.15,0.7);
\end{tikzpicture}
\par\smallskip{\small The willingness to pay of the marginal user, $Vn(1 - n)$, with $V = 10$ and price 1.6.
Between the two equilibria use grows; below the \omterm{def:b3:platforms-and-innovation:network}{critical mass} it collapses.}
\end{omfigure}

The model explains why new networks give access away: a platform that cannot reach
its \omterm{def:b3:platforms-and-innovation:network}{critical mass} at a positive price must subsidise early users until it does.

\section{Pricing on two sides}\label{sec:b3:platforms-and-innovation:two-sided}

\begin{definition}[Two-sided platform]\label{def:b3:platforms-and-innovation:two-sided}
\index{two-sided platform}
A \emph{two-sided platform} serves two groups of users who interact through it, where
each group's value depends on the participation of the other: players and game
developers, card holders and merchants, riders and drivers.
\end{definition}

\begin{proposition}[Two-sided pricing]\label{prop:b3:platforms-and-innovation:two-sided}
Let a monopoly platform face participation $n_i = a_i + \alpha_i n_j - p_i$ on side $i$, where $j$ is the other side, at cost $f_i$
per member. At the profit-maximising prices, $p_i = f_i + n_i - \alpha_j n_j$: each side pays its cost plus a
markup, less the value its members bring to the other side.
\end{proposition}

\begin{proof}
Choose participation instead of prices: $p_i = a_i + \alpha_i n_j - n_i$, so profit is $\sum_i (a_i + \alpha_i n_j - n_i - f_i)n_i$. Its
derivative in $n_1$ is $a_1 + \alpha_1 n_2 - 2n_1 - f_1 + \alpha_2 n_2 = 0$. Since $p_1 = a_1 + \alpha_1 n_2 - n_1$, the condition reads $p_1 - n_1 - f_1 + \alpha_2 n_2 = 0$; side 2
is symmetric.
\end{proof}

The side whose members the other side values most is charged least, and may be
paid to join: its price can fall below cost, or below zero.

\begin{model}{Two-sided platforms}
\assumes A platform sets the price on each side; each side values the other's
participation; users join one platform.
\predicts Prices on each side depend on the other side's valuation; one side can pay
less than its cost.
\tested The markets the model describes are many --- game consoles with players and
developers, payment cards with consumers and merchants, search and advertiser-funded
media, where consumers may dislike the other side --- and a survey of them finds
pricing strategies the model explains alongside outcomes it finds puzzling. % ledger: b3pi:rysman
The model's premise has been tested in one market: estimating consumers' use of
Yellow Pages directories, advertisers' demand and publishers' pricing jointly,
Rysman (2004) found that each side values the other's participation, as the
model assumes. % ledger: b3pi:rys04
Its pricing prediction, that the side valued more pays less, has been described
across markets more than tested.
\breaks When users join several platforms at once (multi-homing), competition moves
to the side that does not, and the price structure changes.
\end{model}

\section{Switching costs}\label{sec:b3:platforms-and-innovation:switching}

\begin{definition}[Switching cost]\label{def:b3:platforms-and-innovation:switching}
\index{switching cost}
A \emph{switching cost} is a cost a buyer bears to move from one supplier to another:
learning a new system, transferring data, breaking a contract.
\end{definition}

\begin{proposition}[Bargains, then rip-offs]\label{prop:b3:platforms-and-innovation:switching}
Two firms with marginal cost $c$ sell a homogeneous good over two periods to buyers who buy
once in each period and bear a \omterm{def:b3:platforms-and-innovation:switching}{switching cost} $s$ to change supplier in the second. In
equilibrium each firm charges its own customers $c + s$ in the second period, and the first-period
price is $c - \beta s$.
\end{proposition}

\begin{proof}
In the second period a rival cannot profitably price below $c$, so a firm can charge its
customers up to $c + s$ without losing them, and charges exactly that. A first-period customer is
therefore worth $p_1 - c + \beta s$ to the firm that wins her. Competition for her drives this to zero, as in
Bertrand: $p_1 = c - \beta s$.
\end{proof}

\omterm{def:b3:platforms-and-innovation:switching}{Switching costs} do not raise lifetime profit here; they shift it, from rip-offs
later to bargains now. The bargain is what hands the profit back, so \omterm{def:b3:platforms-and-innovation:switching}{switching costs}
raise profit only when competition for new buyers cannot hand it back: when firms
cannot price below cost to win them, or when new buyers keep arriving and a firm must
set one price for the customers it has locked in and those it is courting. Buyers who
do not foresee the rip-off are not enough: firms still compete for them, as
\cref{exo:b3:platforms-and-innovation:6} shows.

\section{Patents: the incentive and the deadweight loss}\label{sec:b3:platforms-and-innovation:patents}

\begin{definition}[Patent; patent race; optimal patent length]\label{def:b3:platforms-and-innovation:patent}
\index{patent}\index{patent race}\index{optimal patent length}
A \emph{patent} gives its holder the exclusive right to use an invention for a fixed
term, in exchange for publishing it; its \emph{breadth} is the range of products it
covers. In a \emph{patent race} firms compete to make the same invention first,
since only the first is rewarded. The \emph{optimal patent length} balances the
incentive to invent, which rises with the term, against the deadweight loss of
monopoly during it.
\end{definition}

A longer term raises the reward and so the research effort; it also prolongs the
monopoly. Exercise 11 solves the trade-off in a small model. A \omterm{def:b3:platforms-and-innovation:patent}{patent race} adds a
second distortion: each firm's chance of winning takes rivals' prizes, so firms
together can spend more on the race than the invention is worth to society.

\begin{definition}[Replacement effect]\label{def:b3:platforms-and-innovation:replacement}
\index{replacement effect}
The \emph{replacement effect} is the smaller gain an incumbent monopolist draws from an
innovation, compared with an entrant, because it replaces its own profit.
\end{definition}

\begin{proposition}[The replacement effect]\label{prop:b3:platforms-and-innovation:arrow}
In an industry of price-taking firms with cost $c$, an innovation lowers cost to $c' < c$ without
making the monopoly price at $c'$ fall below $c$. The innovation is worth more to an entrant
holding it than to a monopolist who already serves the market at cost $c$.
\end{proposition}

\begin{proof}
The entrant prices just below $c$ and earns $(c - c')Q(c)$. The monopolist's profit $\pi_m(x)$ at cost $x$ has $\pi_m'(x) = -Q_m(x)$ by
the envelope theorem, so its gain is $\int_{c'}^{c} Q_m(x)\,dx$. For $x \le c$, the monopoly price at $x$ is at least the one at $c'$,
which is at least $c$, so $Q_m(x) \le Q(c)$; the gain is at most $(c - c')Q(c)$.
\end{proof}

\begin{record}[Intellectual property and follow-on research]
The human genome was sequenced twice: by a public project and by a private firm,
whose rights on the genes it sequenced first lasted up to two years, until the
public effort re-sequenced them. Genes held under the firm's rights saw 20 to 30~\% less
later research and product development. % ledger: b3pi:williams
Gene \emph{\omterm{def:b3:platforms-and-innovation:patent}{patents}} gave a different answer. Comparing genes in \omterm{def:b3:platforms-and-innovation:patent}{patent} applications
that were granted with those that were rejected, and using the leniency of the
examiner assigned as an instrument, two designs found no quantitatively important
effect on follow-on research or investment. Patented genes had looked more valuable
before they were patented, which is why a simple comparison misleads. % ledger: b3pi:sw
The two studies differ in the right studied --- a contract-based restriction on
data in one, a published \omterm{def:b3:platforms-and-innovation:patent}{patent} in the other --- and their difference is part of
the evidence.
\end{record}

\begin{omchart}{Patent applications by residents per million people against GDP
per person (PPP, constant 2021 dollars), 2019, 111 countries, both on logarithmic
scales.}{World Bank, World Development Indicators, IP.PAT.RESD, SP.POP.TOTL and
NY.GDP.PCAP.PP.KD}
\begin{tikzpicture}
\begin{axis}[omscatter, width=0.6\linewidth, height=0.4\linewidth,
  xlabel={GDP per person, dollars}, ylabel={Patents per million},
  xtick={3,4,5}, xticklabels={1{,}000, 10{,}000, 100{,}000},
  ytick={-2,-1,0,1,2,3}, yticklabels={0.01, 0.1, 1, 10, 100, 1{,}000}]
  \addplot table[x=lx, y=ly] {figdata/bachelor-3/29-platforms-and-innovation/out/patents.csv};
  \addplot[omThm, thick, no marks, sharp plot] table[x=lx, y=ly] {figdata/bachelor-3/29-platforms-and-innovation/out/patents_fit.csv};
\end{axis}
\end{tikzpicture}
\end{omchart}
% figdata: patents_stats

Patenting rises steeply with income: across countries, 1~\% more income per person goes
with 1.56~\% more resident \omterm{def:b3:platforms-and-innovation:patent}{patent} applications per head, and income accounts for 48~\% of
the variation. The relation says nothing about cause; exercise 8 asks how much research
spending adds.

\section{Competition and innovation}\label{sec:b3:platforms-and-innovation:competition}

\begin{model}{Competition and innovation, the inverted U}
\assumes Firms in an industry are either neck and neck or leader and laggard;
innovation lets a firm escape competition or catch up.
\predicts Competition raises the gain from innovating for neck-and-neck firms and
lowers it for laggards, so innovation first rises and then falls with competition.
\tested U.K. firms' \omterm{def:b3:platforms-and-innovation:patent}{patents} at the U.S. \omterm{def:b3:platforms-and-innovation:patent}{patent} office follow the predicted inverted U
and its other predictions. % ledger: b3pi:abbgh
\breaks The shape depends on how competition is measured --- an inverse of the
price--cost margin is itself shaped by innovation --- and on which firms are
pooled.
\end{model}

\section{Entry, exit and productivity growth}\label{sec:b3:platforms-and-innovation:reallocation}

\begin{definition}[Reallocation decomposition; business dynamism]\label{def:b3:platforms-and-innovation:reallocation}
\index{reallocation decomposition}\index{business dynamism}
A \emph{reallocation decomposition} splits the change in an industry's
share-weighted productivity into the growth of continuing firms (within), the shift
of shares among them (between and cross), and the entry and exit of firms.
\emph{Business dynamism} is the pace of that turnover, measured by entry and exit
rates of firms and of jobs.
\end{definition}

\begin{omchart}{Enterprise birth and death rates, business economy, 32 European
countries, 2022 (percent of active enterprises; Eurostat country codes; values
marked provisional left out).}{Eurostat, business demography (bd\_size), B--S except
public administration and holding companies}
\begin{tikzpicture}
\begin{axis}[omrecord, width=0.62\linewidth, height=0.78\linewidth, xmin=0, xmax=20,
  ymin=-0.8, ymax=31.8, ytick=data,
  yticklabels from table={figdata/bachelor-3/29-platforms-and-innovation/out/dynamism.csv}{geo},
  y tick label style={font=\scriptsize}, xlabel={Percent of active enterprises},
  grid=none, xmajorgrids, legend columns=2,
  legend style={font=\footnotesize, at={(0.5,1.02)}, anchor=south}]
  \addplot[only marks, mark=*, omDef] table[x=birth, y=pos] {figdata/bachelor-3/29-platforms-and-innovation/out/dynamism.csv};
  \addplot[only marks, mark=o, omThm, thick] table[x=death, y=pos] {figdata/bachelor-3/29-platforms-and-innovation/out/dynamism.csv};
  \legend{births, deaths}
\end{axis}
\end{tikzpicture}
\end{omchart}

% figdata: dynamism, dynamism_stats
Across these countries, between 6 and 18~\% of enterprises were born in 2022 (median
10.9~\%) and a median of 8.0~\% died; countries with more births also had more deaths
(correlation 0.57), so high entry goes with high turnover, not only with growth. Most
of the churn is among the smallest firms: in the European Union in 2023, the 10.5~\%
of enterprises that were born employed 2.3~\% of all workers, about one person each,
and the 8.1~\% that died employed 1.7~\%. Four out of five newborn enterprises were
still active a year later and two out of three after two years; those that survived
five years had raised their employment by 77~\% since birth. Entry is thus mostly a
stream of very small trials, of which a few grow; the decomposition below asks how
much of productivity growth they carry.

\begin{proposition}[The FHK decomposition]\label{prop:b3:platforms-and-innovation:fhk}
With shares $s_{it}$ and productivities $p_{it}$, industry productivity $P_t = \sum_i s_{it}p_{it}$ changes by
\[
\Delta P = \sum_{C} s_{i0}\Delta p_i + \sum_{C}\Delta s_i (p_{i0} - P_0) + \sum_{C}\Delta s_i\Delta p_i
  + \sum_{N} s_{i1}(p_{i1} - P_0) - \sum_{X} s_{i0}(p_{i0} - P_0),
\]
where $C$, $N$ and $X$ are continuing, entering and exiting firms.
\end{proposition}

\begin{proof}
Write $P_1 - P_0 = \sum_{C\cup N} s_{i1}p_{i1} - \sum_{C\cup X} s_{i0}p_{i0}$. Since the shares in each year sum to one, $P_0 = P_0\sum_{C\cup N}s_{i1} = P_0\sum_{C\cup X}s_{i0}$,
so $P_1 - P_0 = \sum_{C\cup N}s_{i1}(p_{i1} - P_0) - \sum_{C\cup X}s_{i0}(p_{i0} - P_0)$. For continuing firms, $s_{i1}(p_{i1} - P_0) - s_{i0}(p_{i0} - P_0)$ expands
into $s_{i0}\Delta p_i + \Delta s_i(p_{i0} - P_0) + \Delta s_i \Delta p_i$.
\end{proof}

\begin{measure}[How productivity growth is decomposed]
The decomposition needs establishment data that follow units over time, a
productivity measure (output per worker, or total factor productivity), shares
(of employment or output) and a horizon. Each choice moves the answer: over longer
horizons entrants have had time to learn and exiters to decline, so net entry counts
for more; and the result is cyclical and sensitive to the method. % ledger: b3pi:fhk
Reallocation to more productive units --- among continuing establishments and
through entry and exit --- accounts for a significant share of U.S. manufacturing
productivity growth, and in the services studied net entry dominates.
\end{measure}

\begin{record}[Entry, exit and productivity]
In U.S. retail trade in the 1990s, virtually all productivity growth came from more
productive new establishments displacing much less productive ones that closed. % ledger: b3pi:fhk.retail
In U.K. manufacturing from 1980 to 1992, entry, exit and changes in market shares accounted for
50~\% of establishment labour-productivity growth and 80 to 90~\% of total factor productivity
growth, much of it by firms with several plants closing their poor ones and opening
better ones. % ledger: b3pi:dhh
\end{record}

\section{Exercises}\label{sec:b3:platforms-and-innovation:exercises}

\begin{exercise}[$\star$]\label{exo:b3:platforms-and-innovation:1}
From the opening chart, by how many points of GDP did Korea's research spending rise
from 1996 to 2023? By what factor did China's?
\end{exercise}

\begin{exercise}[$\star$]\label{exo:b3:platforms-and-innovation:2}
Three firms have employment shares 0.5, 0.3 and 0.2 and productivities 10, 20 and 5. The third exits;
the first's productivity rises to 11; an entrant takes a share of 0.2 at productivity 30, and the
others keep their shares. Decompose the change.
\end{exercise}

\begin{exercise}[$\star$]\label{exo:b3:platforms-and-innovation:3}
With $V = 8$ and price 1.5, find the equilibrium shares of users and the \omterm{def:b3:platforms-and-innovation:network}{critical mass}.
\end{exercise}

\begin{exercise}[$\star$]\label{exo:b3:platforms-and-innovation:4}
Two platforms compete in a market said to tip. What data would contradict \omterm{def:b3:platforms-and-innovation:network}{tipping}?
\end{exercise}

\begin{exercise}[$\star\star$]\label{exo:b3:platforms-and-innovation:5}
A platform faces $n_1 = 3 + n_2 - p_1$ and $n_2 = 3 - p_2$ (side 2 does not care about side 1), with costs of 1 on each
side. Find the profit-maximising participation and prices.
\end{exercise}

\begin{exercise}[$\star\star$]\label{exo:b3:platforms-and-innovation:6}
In the switching-cost model, $c = 10$, $s = 4$ and $\beta = 0.9$. Find the prices in both periods and the lifetime
profit per customer. What changes if buyers ignore the second period?
\end{exercise}

\begin{exercise}[$\star\star$]\label{exo:b3:platforms-and-innovation:7}
Demand is $Q = 100 - P$ and the industry's cost is 40; an innovation lowers it to 20. Compare the
gain of a monopolist with that of an entrant.
\end{exercise}

\begin{exercise}[$\star\star$]\label{exo:b3:platforms-and-innovation:8}
Regress the log of resident \omterm{def:b3:platforms-and-innovation:patent}{patent} applications per million people on research
spending as a share of GDP across countries in 2019, then add log income. What happens to
the research coefficient?
\end{exercise}

\begin{exercise}[$\star\star$]\label{exo:b3:platforms-and-innovation:9}
What shape of the relation between competition and innovation in a panel of
industries would refute the inverted U?
\end{exercise}

\begin{exercise}[$\star\star\star$]\label{exo:b3:platforms-and-innovation:10}
Prove the two-sided pricing rule, and show that a side can be paid to join.
\end{exercise}

\begin{exercise}[$\star\star\star$]\label{exo:b3:platforms-and-innovation:11}
An invention saves $\sqrt R$ million dalens a year if $R$ is spent on research. While patented, its holder
earns half the saving and the other half is lost; after expiry users get it all. With
discount rate 5~\% and $A(T) = (1 - e^{-0.05T})/0.05$, the inventor chooses $R$ to maximise $\tfrac12\sqrt R\,A(T) - R$. Find the
\omterm{def:b3:platforms-and-innovation:patent}{patent} length that maximises welfare.
\end{exercise}

\begin{exercise}[$\star\star\star$]\label{exo:b3:platforms-and-innovation:12}
On an invented two-year panel built from Brenland's manufacturing firms, compute the
decomposition of the change in output per worker.
\end{exercise}

\section{Problem: Brenmoor's Ride Platform}\label{sec:b3:platforms-and-innovation:problem}

\begin{problem}[{Weekend problem --- Brenmoor's ride platform}]\label{pb:b3:platforms-and-innovation:ride}
A ride platform starts in Brenmoor. % figdata: ride_stats

\medskip\noindent\textbf{Part I --- reaching riders.} Of Brenmoor's 10{,}000 residents, one
with taste $\theta$ (uniform on $[0, 1]$) values a month of the service at $10\theta n$ dalens when a share $n$ of residents
uses it. The monthly price is 1.6.
\begin{enumerate}
\item Find the equilibrium shares of users.
\item Which is the \omterm{def:b3:platforms-and-innovation:network}{critical mass}, in residents?
\item The platform launches with 1{,}000 users. What happens?
\item What is the highest price at which any positive share can be an equilibrium?
\item How many free months of service would you need to give, and to how many, to
pass the \omterm{def:b3:platforms-and-innovation:network}{critical mass}?
\item Why is the share 0 an equilibrium at any positive price?
\end{enumerate}

\medskip\noindent\textbf{Part II --- two sides.} With $n_R$ and $n_D$ the riders and drivers (hundreds),
$n_R = 4 + 1.2n_D - p_R$ and $n_D = 2 + 0.4n_R - p_D$, where $p_R$ and $p_D$ are the platform's monthly fees (a negative fee is
a payment); the platform's cost is 1 per member on each side.
\begin{enumerate}[resume]
\item Write the platform's profit as a function of $n_R$ and $n_D$.
\item Derive the two first-order conditions.
\item Solve for $n_R$ and $n_D$.
\item Find the fees.
\item Check the pricing rule of the chapter on each side.
\item Which side is subsidised, and why?
\item Find the platform's profit.
\item If riders valued drivers less, at $\alpha_R = 0.8$, would the drivers' fee rise or fall?
\end{enumerate}

\medskip\noindent\textbf{Part III --- a rival.}
\begin{enumerate}[resume]
\item A rival platform enters, and drivers sign up with both. On which side will the
platforms compete hardest?
\item What does multi-homing by drivers do to the drivers' fee?
\item Why does multi-homing by drivers make \omterm{def:b3:platforms-and-innovation:network}{tipping} less likely?
\item If drivers could not multi-home, how much would a rival have to give each rider, per month,
to offset the incumbent's network of drivers?
\end{enumerate}

\medskip\noindent\textbf{Part IV --- the services sector.} Take Brenland's services firms
in the 2025 census (\code{load("firms")}) as year 0, with output per worker as productivity and
employment shares. Two years later the tenth with the lowest productivity have
exited, continuing firms are 2~\% more productive with the same workers, and 30 entrants
each have the median workforce and productivity 10~\% above the year-0 mean.
\begin{enumerate}[resume]
\item Find industry productivity in both years and its growth.
\item Decompose the change into within, between, cross, entry and exit.
\item What share of the change comes from entry and exit together?
\item Why is the exit term so large though the exiters are only a tenth of firms?
\item Compare with the evidence of the record box.
\item State the result: the platform's fees on each side, and the share of the
sector's productivity growth that comes from entry and exit.
\end{enumerate}
\end{problem}
